\section{Design levers II: charge management and interfacial catalysis}
\label{sec:charge}

\subsection{L3 --- Excitons, crystallinity, and stacking order}
\label{sec:excitons}

Absorbed photons in organic frameworks yield strongly bound excitons, so the third lever is
whether those excitons dissociate and whether the resulting charges survive transport.
Mechanistic perspectives now treat photon absorption, exciton dissociation, and electron
delivery as separately measurable stages of one pipeline~\cite{blatte2024photons}.

\textbf{Excited-state design (L3a).} Computation has moved upstream of synthesis:
density-functional screening across donor--acceptor pairs shows the exciton binding energy
$E_b$ is systematically tunable through D--A strength, making $E_b$ a design variable rather
than a post-hoc measurement~\cite{qian2023computation}. On the experimental side, an
isoreticular hydrazide-COF series resolved how excited-state electron distribution and
transport pathways change with the linker, tying photophysics to specific framework
positions~\cite{chen2023tuning}. What is still missing is the closed loop between the two
communities --- computed $E_b$ values validated by spectroscopy on the same material series
(Sec.~\ref{sec:challenges}, G5)~\cite{qian2023computation,chen2023tuning,blatte2024photons}.

\textbf{Crystallinity and stacking (L3b).} The factor-deconvolution benchmark identified
crystallinity as a prime HER factor~\cite{ghosh2020identification}, and two complementary
studies show why and how to act on it. During photocatalytic cycling, coaxial stacking can
disorder and activity decays; threading poly(ethylene glycol) through the mesopores locks the
stacked configuration and prevents the decline~\cite{zhou2021peg} --- direct evidence that
\emph{operando} structural integrity, not just as-made crystallinity, matters. At the synthesis
end, the reconstruction strategy decouples polymerization from crystallization: amorphous
imine networks pre-organized by a removable covalent tether reconstruct into highly crystalline
frameworks, lifting sacrificial HER to 27.98~mmol\,h$^{-1}$\,g$^{-1}$ --- the highest
sacrificial rate in this survey's library --- while confirming that poor reversibility of
robust linkages is the central obstacle the field must engineer
around~\cite{zhang2022reconstructed,haase2020solving}.

\subsection{L4 --- Cocatalysts and interfacial engineering}
\label{sec:interface}

Even perfectly delivered electrons produce no hydrogen without an interface that turns them
over. The library evidence organizes into three interface strategies.

\textbf{Pt speciation (L4a).} Platinum remains the default electron sink, but its speciation is
now a lever in its own right: dispersing Pt as single atoms on a COF raises active-site density
relative to nanoparticles~\cite{dong2021platinum}, while in-situ photodeposition yields
cluster-scale Pt with strong electronic coupling to the framework~\cite{li2022in}. The
factor-deconvolution protocol treats Pt loading and speciation as one of the deconvoluted
variables, which is what makes cross-material activity comparisons meaningful at
all~\cite{ghosh2020identification}.

\textbf{Molecular and earth-abundant cocatalysts (L4b).} The azine N$_2$-COF paired with a
chloro(pyridine)cobaloxime delivers 782~$\mu$mol\,h$^{-1}$\,g$^{-1}$ at a turnover number of
54.4, with mechanistic analysis supporting outer-sphere electron transfer to a monometallic
cobalt pathway~\cite{banerjee2017single}. Covalently wiring the cobaloxime to the backbone
improves and prolongs activity relative to physisorbed analogues by enabling recoordination of
the cocatalyst during operation~\cite{gottschling2020rational}. These molecular systems remain
orders of magnitude below Pt-based benchmarks in rate and turnover
(cf.~\cite{yang2021protonated}), a gap analyzed in Sec.~\ref{sec:challenges}
(G6)~\cite{banerjee2017single,gottschling2020rational}. Interface design more broadly ---
including cocatalyst anchoring chemistry --- is a central theme of the 2024 Accounts
survey~\cite{chen2024photocatalysis}.

\textbf{The sacrificial-system context.} Almost all quantitative claims in this section are
measured against sacrificial electron donors, and the donor is itself an interface-level
variable: the cobaloxime/N$_2$-COF benchmark operates with triethanolamine in a
water/acetonitrile medium~\cite{banerjee2017single}, while the record-setting protonated D--A
systems were developed in ascorbic-acid media, where the donor also participates in the
protonation equilibrium that red-shifts absorption~\cite{yang2021protonated}. Cross-study
comparisons that ignore the donor identity therefore conflate interfacial kinetics with donor
thermodynamics --- one of the reasons a fixed-protocol benchmark is registered as gap G1 in
Sec.~\ref{sec:challenges}~\cite{ghosh2020identification,banerjee2017single}.

\textbf{Heterojunctions and Z-schemes (L4c).} Coupling COFs to inorganic partners re-routes
charges across a built-in junction: an in-situ grown, chemically bonded 2D/2D COF /
oxygen-vacancy-WO$_3$ Z-scheme targets \emph{overall} water splitting rather than sacrificial
HER~\cite{shen2023in}, and a 2D/2D COF/CdS Z-scheme heterojunction enhances photocatalytic
hydrogen evolution~\cite{gao20222d}. Composite and heterojunction photocatalysts are surveyed
systematically in the 2020 CSR review~\cite{wang2020covalent}. The overall-water-splitting
rates remain far below sacrificial benchmarks~\cite{shen2023in}, marking the sacrificial-donor
dependency as a field-level open problem (Sec.~\ref{sec:challenges}, G2).

Taken together, L3 and L4 close the causal chain begun in Sec.~\ref{sec:levers}: linkage and
band engineering determine what can be absorbed and separated; crystallinity, stacking, and the
catalytic interface determine how much of that potential is converted. The prime-factor
evidence~\cite{ghosh2020identification} justifies the ordering used in the synthesis design of
Sec.~\ref{sec:routec}: crystallinity first, interface second, everything else after.
